\section{End-to-end pipeline architecture}
\label{sec:architecture-pipeline}

This part describes the effective path of a LANCE audit, from the selection of an
execution mode to the files deposited in the run folder, relying exclusively on
the code of the \code{dev/1} branch. The statements refer to the module and,
where useful, to the function that implements them; behaviours that the code does
not allow to be established are flagged as such rather than filled in by
inference.

\subsection{Overview and mode selection}

The application does not carry the audit logic: it validates the request,
guarantees a single active run per instance, streams the SSE events and re-reads
the artefacts (\code{docs/architecture/run-modes.md}). It is \code{make\_runner}
(\code{src/agent/run\_modes.py}) that chooses the execution strategy. For
\code{lance} and \code{rules}, it instantiates a \code{Pipeline} and wraps it in
\code{LanceRunner}, an adapter that reuses the pipeline object and its public
API. For \code{vanilla} and \code{scripted}, it builds a \code{SimpleRunner}, a
standalone class that shares with the pipeline only the tool execution boundary,
the evidence log, the cost tracking and the laboratory reservation.

\begin{table}[H]
\centering
\caption{Execution modes and corresponding boundary.}
\label{tab:run-modes}
\small
\begin{tabular}{@{}llll@{}}
\toprule
Mode & Engine & Model & Effects \\
\midrule
\code{lance} & \code{Pipeline} (6 phases) & required & scenario, deployment, evaluation \\
\code{rules} & \code{Pipeline} (phases 3--4) & none & scenario and deterministic audit \\
\code{vanilla} & \code{SimpleRunner} & required & bounded Nmap probe on CIDR \\
\code{scripted} & \code{SimpleRunner} & none & versioned procedures \\
\bottomrule
\end{tabular}
\end{table}

The \code{rules} mode is not a second engine: \code{make\_runner} sets
\code{decision\_policy="rules"}, \code{phases=[3, 4]} and
\code{execution\_profile="full"}, and rejects any provider, direct network or
monetary budget. The \code{vanilla} and \code{scripted} modes, conversely, reject
\code{scenario\_id} and any harness-specific setting. The qualifier
\code{full}/\code{compact} designates the orchestration of the pipeline, never the
type of runner.

\subsection{The six phases and their deliverables}

The declarative registry \code{src/agent/registry.py} (\code{AGENTS}) assigns to
each phase a prompt template, a tool group, a validator and a deliverable name.
The routing is explicit in \code{src/agent/phases/registry.py}:
\code{PHASE\_MODULES} maps the numbers 1 to 6 to fixed modules, without dynamic
import.

\begin{table}[H]
\centering
\caption{Phases, modules, main deliverables and associated artefacts.}
\label{tab:phases}
\small
\begin{tabular}{@{}lp{2.7cm}p{3.2cm}p{4.05cm}@{}}
\toprule
Phase & Module & Main deliverable & Associated artefacts \\
\midrule
1 & \code{graph/run.py} & \code{01\_graph\_analysis.md} & \code{01\_graph\_evidence.json} \\
2 & \code{recon/run.py} & \code{02\_recon.md} & \code{02\_recon\_evidence.json}, \code{02\_topology\_edges.json} \\
3 & \code{analysis/run.py} & \code{03\_vuln\_analysis.json} & \code{03\_vuln\_analysis\_raw.json}, \code{03\_phase3\_status.json}, \code{03\_device\_*} \\
4 & \code{verification/run.py} & \code{04\_exploitation.json} & \code{execution\_errors} diagnostics \\
5 & \code{intrusion/run.py} & \code{05\_intrusion.json} & \code{05\_intrusion\_context.json} \\
6 & \code{report/run.py} & \code{06\_report.md} & \code{06\_report\_analysis.md}, per-section sheets \\
\bottomrule
\end{tabular}
\end{table}

Phase 1 projects the results of the graph tools into a log of facts independent
of the model's discourse
(\code{GraphPhase.\_build\_graph\_evidence\_projection}), and the runner rejects
a submission whose per-device coverage is incomplete
(\code{src/agent/core/runner.py}). Phase 2 lets the model choose the probes, but
\code{02\_recon.md} is rendered from the tool log, not copied by the model. Phase
3 runs in three steps: \code{scan.scan\_phase} (deterministic observations),
\code{devices.analyze\_devices} (analysts per device) and then
\code{run.aggregate\_phase}, which normalises, deduplicates and projects the
findings and writes \code{03\_vuln\_analysis.json}. Phase 4 launches micro-agents
per vulnerability and aggregates their results. Phase 5 produces its assessment
from the log: the model supplies only an end marker, and the status remains
\code{objectives\_status=not\_certified}. Phase 6 is particular: its
\code{AgentConfig} declares no tool and \code{max\_turns=1}; the report is
assembled deterministically from sheets, the model writing only a bounded analyst
note (\code{06\_report\_analysis.md}).

\subsection{Graph and run context}

The graph is built in two layers. \code{src/loader.py} reads the infrastructure
YAML into dataclasses (\code{load\_yaml}) and then populates a backend
(\code{build\_graph}, default backend \code{NetworkXBackend}).
\code{src/graph\_backend.py} defines the abstract interface \code{GraphBackend}
(\code{get\_device}, \code{get\_neighbors}, \code{find\_all\_paths},
\code{get\_attack\_surface}, \code{get\_graph\_stats}, \code{to\_dict}) and its
NetworkX implementation on a \code{DiGraph}. Two details deserve to be stressed:
\code{get\_neighbors} returns the sorted union of successors and predecessors, and
\code{find\_all\_paths} works on an undirected view (\code{to\_undirected}) in
order to find paths in both directions, including through the abstract backend.

The state belongs to the run. Each \code{Pipeline} creates an immutable
\code{RunContext} (run folder and \code{GraphContext}); \code{GraphContext}
carries the backend, the infrastructure, the CVE reports, the risk scores, the
scenario or discovery topology, the weighted graph and the latest delta. The graph
tools retrieve this context via a \code{ContextVar}
(\code{src/agent/core/graph\_context.py}), without modifying the arguments exposed
to the model; \code{graph\_bound} and \code{bind\_graph\_callable} explicitly
capture the owner, including for thread workers, rather than relying on implicit
context propagation.

\subsection{The tool governance boundary}

All profiles and all modes go through \code{wrap\_tool}
(\code{src/agent/core/executor.py}). On each call, the wrapped function:
(1) checks that the previous piece of evidence could be archived, cooperative
stop and the budgets via \code{check\_execution\_limits}; (2) atomically
increments a counter and refuses the call if \code{max\_tool\_calls} is reached;
(3) assigns a reference \code{tc-<uuid>}; (4) applies the scope refusals;
(5) executes the tool; (6) systematically archives an entry in
\code{tool\_calls.jsonl}, including for a refusal or an exception.

\begin{table}[H]
\centering
\caption{Controls applied by \code{check\_execution\_limits} and \code{wrap\_tool}.}
\label{tab:executor}
\small
\begin{tabular}{@{}p{2.6cm}p{4.6cm}p{6.0cm}@{}}
\toprule
Control & Mechanism & Consequence of a failure \\
\midrule
Evidence integrity & flag \code{\_evidence\_integrity\_failed} & \code{EvidenceWriteError}, run not verifiable \\
User stop & \code{\_stop\_event} (experimental scope) & \code{RunStopped} \\
Time budget & \code{max\_duration\_s} via \code{time.monotonic()} & \code{BudgetExceeded} \\
Monetary budget & \code{tracker.check\_budget()} & \code{BudgetExceeded} \\
Call budget & counter under \code{\_artifact\_log\_lock} & \code{BudgetExceeded} \\
Scope & \code{\_intrusion\_scope\_violation} & refused call archived as JSON \\
Free execution & \code{python\_exec} & refusal \code{unverifiable\_execution\_scope} \\
\bottomrule
\end{tabular}
\end{table}

Two properties are important for the interpretation of the results. On the one
hand, a failure to write the log is never ignored: it is converted into
\code{EvidenceWriteError}, which interrupts the run rather than leaving untraced
actions. On the other hand, the cost limit of \code{check\_budget} applies to
observed usage (\code{total\_cost()}), not to a reservation of in-flight calls;
the code docstring says so explicitly. The scope is checked on the destination
fields (\code{ip}, \code{host}, \code{target}, \code{url}) against the
\code{target\_subnet} of the run; an action without a verifiable destination is
refused upstream (\code{python\_exec}), and the module header specifies that these
are application guards, not a system sandbox.

\subsection{Artefacts and cost tracking}

The run folder is created exclusively under \code{output/agent/}
(\code{Pipeline.\_\_init\_\_}); in the event of a collision, a hexadecimal suffix
is appended and no earlier folder is overwritten. \code{run\_meta.json} records
phase, profile, policies, commit and limits; \code{cost\_summary.json} is written
in the \code{finally} block of the pipeline, and therefore even on exception. The
collection contract (\code{src/benchmark/artifacts.py}) defines a positive list of
allowed files (\code{\_ALLOWED\_ROOT\_FILES}) and \code{03\_device\_*.json}
patterns; \code{create\_submission\_bundle} produces a zip equipped with a
SHA-256 manifest and re-checks the manifest after writing. On the agent side, a
deliverable attempt is archived in \code{.attempts/} and recorded in
\code{deliverable\_attempts.jsonl}; the route
\code{GET /api/runs/\{id\}/download/zip} excludes private files and hidden
entries.

Cost tracking (\code{src/agent/cost\_tracker.py}) resolves the price once per
phase (\code{\_resolve\_pricing}) and records it in \code{PhaseUsage}, so that a
completed run cannot be re-evaluated after the fact. Each phase carries the source
of the price (\code{dynamic\_catalog}, \code{static\_catalog},
\code{subscription}, \code{default\_estimate}) and a \code{cost\_is\_estimate}
indicator. \code{summary()} explicitly distinguishes the cumulative duration of
the agents (\code{total\_agent\_duration\_s}) from the wall-clock time
(\code{wall\_clock\_duration\_s}); the former does not measure the time actually
elapsed.

\subsection{What the architecture does not guarantee}

The architecture is isolated per run for the deliverables and the graph, but it is
not entirely so: \code{src/agent/phases/README.md} indicates that the phase
services may still call the common engine and that the simultaneity of scenarios
on a shared laboratory is not guaranteed. \code{reserve\_lab} is only active if
\code{LANCE\_LAB\_LOCK} is set and is only reentrant in the same thread. Moreover,
without an explicit parameter, the interactive pipeline imposes neither a call
ceiling nor an analysis timeout per device: the budgets of
table~\ref{tab:executor} are therefore conditional.

Three points could not be established by reading the code. \emph{One}: no
mechanism for archiving or compacting run folders was found in \code{src/} (no
occurrence of \code{tarfile}, \code{shutil.make\_archive} or equivalent);
retention appears to be that of the file system, which the documentation does not
describe. \emph{Two}: \code{02\_recon\_evidence.json} is read by
\code{src/agent/scanner.py} and \code{src/agent/core/runner.py}, but its main
producer has not been identified; asserting it would be a supposition.
\emph{Three}: \code{provider\_events.jsonl} is recorded by the provider callback
(\code{core/runner.py}) via \code{provider\_diagnostics}, and the exact content of
the fields has not been verified line by line. Finally, a \code{completed} status
means that the runner has finished its work, never that the target is exhaustively
covered nor that an access has been proven
(\code{docs/architecture/run-modes.md}).
